\orientacao{Máximo de 15 páginas úteis, excluindo capa, resumo, palavras-chave, declaração de utilização de IA e ferramentas de apoio, índice e referências bibliográficas. Sem limite de referências bibliográficas. Não são admitidos anexos ou apêndices. Respeitar a formatação do template (Arimo (metricamente compatível com Arial), 10 pt, espaçamento simples e margens de 25 mm). Os textos orientadores estão em itálico e a cinzento. O texto Lorem ipsum exemplifica a formatação do corpo do documento (preto, sem itálico). Eliminar ambos na versão de entrega.}
\orientacao{A estrutura apresentada neste template é meramente indicativa. O estudante deverá definir a organização do documento, incluindo a designação, o número e a sequência das secções e subsecções, em função da natureza, da especificidade e dos objetivos da Dissertação, Projeto ou Estágio.}
\orientacao{Independentemente da estrutura adotada, o documento deverá contemplar a contextualização do tema, a revisão e análise crítica do conhecimento relevante e o posicionamento do trabalho a desenvolver. As figuras e tabelas são meramente exemplificativas, tanto na estrutura como no conteúdo, e não são de inclusão obrigatória.}

\section{Introdução e âmbito da revisão}
\orientacao{Contextualizar o tema da D/P/E e delimitar o âmbito da revisão técnico-científica. Identificar os principais aspetos a aprofundar e justificar a sua relevância para o trabalho a desenvolver.}

Lorem ipsum dolor sit amet, consectetur adipiscing elit. Sed do eiusmod tempor incididunt ut labore et dolore magna aliqua. Ut enim ad minim veniam, quis nostrud exercitation ullamco laboris nisi ut aliquip ex ea commodo consequat. Duis aute irure dolor in reprehenderit in voluptate velit esse cillum dolore eu fugiat nulla pariatur.